\chapter{Die Cincinnatier Angelegenheit}
\quellseite{2}

Ueber diese Angelegenheit hat man in unserer Synode im Laufe der letzten zwei Jahre durch Veröffentlichungen verschiedener Art Seitens der dortigen Gemeinde mancherlei gehört, wodurch die Gemüther mancher beunruhigt worden sind. In den betreffenden Veröffentlichungen und Circularen, die auch zum Theil in die öffentlichen Zeitungen übergegangen sind, ist wiederholt die Handlungsweise der Synodalbeamten, welche von Amts wegen in dem Falle zu handeln hatten, als eine falsche verdächtigt worden,
\quellseite{3}
und man möchte nun gern — so wird uns von vielen Seiten versichert — von den Beamten eine öffentliche Erklärung hören. Zwar liegt die ganze Angelegenheit zur Zeit in den Händen des Mittleren Districts, dessen Präses sich genöthigt gesehen hat, die beiden Pastoren der Gemeinde bis zur nächsten Synode von der Synodalgemeinschaft zu suspendiren, und der seinem District darüber berichten und dann in der Lage sein wird, sich gegen die wider ihn erhobenen Beschuldigungen zu rechtfertigen. Anders steht es in Bezug auf die Beamten der Allgemeinen Synode. Die Allgemeine Synode versammelt sich erst im Jahre 1905. Sodann haben öffentliche Zeitungen, sowohl weltliche als kirchliche, die gegen die Beamten der Synode erhobenen Beschuldigungen verbreitet.\footnote{Es ist sehr zu bedauern, daß solche Dinge, die noch in den geschlossenen Kreis einer Synode gehören, in die Oeffentlichkeit gezerrt werden. Schon \gesperrt{vor} aller Untersuchung Seitens der Districtscommission und natürlich auch \gesperrt{vor} der Suspension Seitens des Districtspräses erschien in einer Cincinnatier Zeitung (vom 7. September 1903) das folgende „Eingesandt“: „Von der Kanzel der ev.-luth. Dreifaltigkeits-Kirche wurde der Gemeinde von Herrn P. Alexander v. Schlichten gestern Folgendes bekannt gegeben: ‚Die Ehrw. Commission, welche von unserer Districtssynode in Cleveland ernannt worden ist, bestehend aus Herrn Vicepräses P. Dau aus Hammond, Herrn Prof. Zucker aus Fort Wayne und dem Laiendelegaten Herrn Ahlbrandt aus Seymour, Ind., hat angezeigt, daß morgen Abend über 8 Tage, also am 14. September, eine Gemeindeversammlung abgehalten werden soll. Im Auftrag der Synode soll diese Commission zwei Dinge thun: 1. den Bannfall A.’s untersuchen, da derselbe von dem Allgemeinen Präses als ein falscher Bann erklärt worden ist; und 2. soll die Commission das rechte Verhältniß zwischen unserer Gemeinde und der sogenannten Savior Church herzustellen suchen. Gewissenshalber muß ich der Gemeinde nun noch folgende Erklärung abgeben: Ich habe bis jetzt während des ganzen langen Kampfes immer und allezeit die größtmöglichste Rücksicht genommen, sowohl auf unsere Synodalbeamten, sowie auch auf die kleine Opposition innerhalb unserer Gemeinde. ... Jetzt aber nach der Synode in Cleveland, woselbst ich genug und abergenug frei und offen geredet habe, ist’s endlich mit diesem ewigen Rücksichtnehmen auf immer aus und vorbei. Nun werde ich mit Gottes Hülfe ganz auf dem Plan sein und beweisen, daß wir uns vor keinem Menschen auf der Welt zu fürchten brauchen: weder vor unsern Beamten, noch vor der englischen Synode, noch vor jener sogenannten Savior Church, früher “Christ Church”, noch vor irgend welcher Opposition in unserer Gemeinde. Die Wahrheit soll und muß endlich triumphiren! — Die berühmten Worte des größten deutschen Staatsmannes, des Fürsten Otto v. Bismarck, der ganzen Welt zugerufen im Reichstag im Februar 1888, eignen wir uns jetzt an in folgender Form: „Die deutsche ev.-luth. Dreifaltigkeits-Gemeinde in Cincinnati, Ohio, und deren Pastoren und Schullehrer fürchten Gott und sonst nichts auf der Welt!“‘“ Ebenso erfahren wir auch bereits aus \gesperrt{Zeitungen}, daß der Theil der Gemeinde, welcher das Circular an die „Glaubensbrüder nah und fern“ gesandt hat, \gesperrt{aus der Synode ausgetreten ist.}} So soll denn auf vielseitiges Begehren über die Betheiligung des Allgemeinen Präsidiums an dem Handel hier Folgendes zur allgemeinen Kenntniß gebracht werden.

Wir knüpfen dabei an das letzte Circular der Cincinnatier Gemeinde vom 13. December 1903 an, das man wohl ihr Ultimatum nennen kann. In demselben heißt es: „Was nun die einzelnen Punkte der amtlichen Anzeige“ (die Suspensionserklärung ist gemeint) „betrifft, so erklären wir zum \gesperrt{ersten} Punkt: Noch niemand, weder Herr Präses Niemann, noch der Hochw. Allgemeine Präses ... hat uns aus Gottes Wort nachzuweisen vermocht, daß jener Selbstbann ein \gesperrt{falscher} sei.“ — Hierzu ist Folgendes zu bemerken: Im Sommer 1902 erklärte die Gemeinde in Cincinnati in einem Circular an die Gemeinden der Synodalconferenz, daß sie die Evangelical Lutheran Church of Our Savior ihres Ortes, die mit der Synodalconferenz verbunden war, nicht als Schwestergemeinde anerkenne, und dies darum, weil jene Organisation theils aus Personen bestehe, mit denen sie bereits seit 22 Monaten im Kampfe stehe, theils aus Personen, die in einer unordentlichen, Aergerniß erregenden, schriftwidrigen Weise, nämlich durch einfache Zusendung eines Resignationsschreibens, ihrer alten Gemeinde den Rücken gekehrt hätten. Dieses Circular wurde für den Präses der englischen Synode Veranlassung, zur Zeit der Synodalconferenz bei unserm Allgemeinen Präses eine Klage einzureichen und eine Untersuchung der Sachlage zu beantragen. Dieser trat dann mit der Gemeinde in Cincinnati in Unterhandlung, und endlich, Mitte October 1902, war dieselbe zu einer Untersuchung bereit. Bei derselben waren außer dem Allgemeinen Präses auch die beiden Vicepräsides zugegen. Um beurtheilen zu können, ob wirklich, wie das Circular besagte, die englische Gemeinde sich aus Leuten zusammensetze, welche eigentlich Glieder der deutschen Gemeinde seien und von derselben keine ordentliche Entlassung erhalten hätten, baten wir die Gemeinde, uns die Namen der betreffenden Personen zu nennen und anzugeben, in welchem Verhältniß dieselben zu ihr gestanden hätten. Anfangs weigerte man sich entschieden, dieser Bitte zu willfahren. Immer wieder wurde uns gesagt, wir brauchten diese Namen nicht zu wissen. Es gehe den Präses nichts an, wer die Leute seien und was die Gemeinde mit denselben gehandelt habe. Der Präses habe kein Recht, sich in ihre Angelegenheiten zu mischen. Wir stellten der Gemeinde vor, daß wir ohne genaue Kenntniß der einschlägigen Thatsachen unmöglich in dem Handel urtheilen könnten; die Gemeinde habe doch auch selbst durch ihr Circular an die Gemeinden der Synodalconferenz zu solcher Untersuchung aufgefordert. Die Antwort war: jenes Circular sei gar nicht als eine solche Aufforderung gemeint gewesen. Der Vorstand sonderlich erklärte, die Gemeinde habe mit jenem Schreiben nur den auswärtigen Brüdern zur Kenntniß bringen wollen, was von der Church of Our Savior zu halten sei, nach 2 Thess. 3, 14. Die Gemeinde in Cincinnati habe ihr Urtheil über die sogenannte englische Gemeinde gefällt. Dieses Urtheil habe sie der Kirche zur Kenntniß gebracht, und jeder Christ habe dieses Urtheil ohne Weiteres anzunehmen und demgemäß zu handeln. Und dies nach dem Wort des HErrn: „Höret er die Gemeine nicht“ etc. Immer wieder sagte man uns, wir sollten uns über das Princip aussprechen, ob eine von einem \gesperrt{Selbstgebannten} gestiftete Gemeinde eine rechte christliche Gemeinde sein könne. Da sei das Protokoll der Gemeinde. Darnach sollten wir urtheilen. Endlich jedoch, nach langem Reden und Kämpfen, erhielten wir die gewünschte Auskunft, und damit vertagte sich die Versammlung bis zum folgenden Abend. Wir sagten uns nun aber, daß wir jetzt nicht umhin könnten, den vielerwähnten Selbstbann des Herrn A. nach dem Protokoll der Gemeinde zu prüfen. Dies thaten wir am nächsten Tag und fanden Folgendes: Herr A. hatte seinen Sohn, nachdem derselbe nicht ganz bis zum zwölften Jahre die Gemeindeschule besucht hatte, aus der Gemeindeschule genommen und in die Freischule geschickt, weil sein Sohn in den weltlichen Fächern in der Gemeindeschule nicht genug lerne. Später werde er denselben in den Confirmandenunterricht schicken. Außerdem hatte Herr A. in der Gemeindeversammlung eine Schrift vorgelegt, in welcher er die Zweckmäßigkeit des Instituts der Gemeindeschule unter den dortigen Verhältnissen leugnete und ein Ersatzmittel vorschlug, durch welches die Gemeinde nach seiner Meinung ihre
\quellseite{4}
Pflicht gegen die heranwachsende Jugend besser erfüllen könnte. Darüber wurde er von der Gemeinde in Zucht genommen. Er zog die Schrift zurück, erklärte, daß er darin zu weit gegangen sei und sich versündigt habe, und bat die Gemeinde, ihm zu vergeben. Er versprach auch im Laufe der Verhandlungen, seinen Sohn im September nächsten Jahres, dann schon mit Neujahr wieder in die Schule senden zu wollen; um der Gemeinde bekannten Umstände willen möchte man bis dahin Geduld haben. Die Gemeinde bestand jedoch darauf, daß A. seinen Sohn schon am folgenden Tag schicken müsse. Man setzte ihm sehr zu, und es fehlte dabei auch nicht an wirklich beleidigenden Aeußerungen Seitens des Pastors. A. verließ dann die Versammlung, „nicht im Zorn, sondern ruhig und langsam“, wie es in der sehr ausführlichen Bann-Erklärung heißt, und die Gemeinde faßte den Beschluß, nachdem einige Glieder noch vergeblich um Aufschub gebeten hatten, Herrn A. als einen solchen anzusehen, der sich selbst in den Bann gethan habe. — Wir konnten nicht anders urtheilen, als daß diese Bann-Erklärung eine ungerechtfertigte war. Herr A. hatte in der Hauptsache die Gemeinde gehört, war auf dem besten Wege, wieder ganz zurechtzukommen. Er bat, was die Sendung seines Sohnes anlangte, um etwas Geduld häuslicher Verhältnisse wegen, worin auch die sachliche Erklärung lag, noch weiter mit sich handeln lassen zu wollen. Die Gemeinde konnte also unmöglich erkennen, daß A. die Gemeinde nicht höre, und folglich konnte sie ihn auch nicht für einen „Selbstgebannten“ erklären, der nach Matth. 18 für einen Heiden und Zöllner zu halten sei. Was sollten die Beamten der Allgemeinen Synode nun thun? Bannfälle zu untersuchen ist nicht eigentlich Sache des Allgemeinen Präsidiums, kann es nur auf dem Appellationswege werden. Wir waren auch gar nicht zu dem Zweck nach Cincinnati gekommen. Und doch war uns klar, wir durften um der Liebe willen nicht verschweigen, was wir gefunden hatten, durften es um so weniger, als die Gemeinde selbst uns genöthigt hatte, uns mit des Herrn A. Sache zu beschäftigen, wenn auch die Gemeinde nicht daran dachte, daß wir ihr Kirchenzuchtverfahren beurtheilen sollten. Wir beschlossen daher, uns in dem der Gemeinde zu unterbreitenden Gutachten auch über diese Sache auszusprechen. Daß dies in möglichst schonender Weise geschehen ist, wird der Leser selbst erkennen. Folgendes ist nämlich der Wortlaut des Gutachtens, welches wir am zweiten Abend, am 17. October 1902, der Gemeinde vorlegten:

\section{Gutachten}

\subsection{I}

Wir erkennen mit großer Freude, wie in allen Verhandlungen der Dreifaltigkeits-Gemeinde die Sorge für die Pflege und Erhaltung ihrer Gemeindeschule zum Ausdruck kommt. Die Gemeinde hat große Opfer gebracht für die so wichtige Schulsache. Sie hat wohl daran gethan, zu jeder Zeit Leuten die Stirne zu bieten, welche den Versuch machten, dies wichtige, weil für die Gemeinde so nöthige und segensreiche Werk zu hindern und zu zerstören. Hierin stand und steht sie auf dem Grund und im Gehorsam des Wortes Gottes. Gott erhalte sie unentwegt in dieser Stellung, nämlich in treuer Pflege ihrer Schule und, wo sie dazu gezwungen wird, im heiligen Kampf um dies theure Gemeindekleinod. Es ist dies ja unsere Stellung, die Stellung unserer Synode, unserer Gemeinden, wie denn je und je von uns bezeugt worden ist, daß es zur rechten Gestalt einer evangelisch-lutherischen Ortsgemeinde gehört, daß sie christliche Schulen einrichtet und pflegt.

Doch dürfen wir der werthen Dreifaltigkeits-Gemeinde auch nicht verhehlen, daß wir bei Durchlesung der Protokolle über ihren löblichen und wohlberechtigten Kampf, den die Gemeinde für ihre Schule in treuer Sorge geführt hat, auf einen Punkt gestoßen sind, über den wir der Gemeinde eine Erklärung geben müssen, die wir ihr gewiß nicht deshalb verschweigen oder vorenthalten dürfen. Wir meinen nämlich den Fall des Herrn A., den die Gemeinde für einen Selbstgebannten erklärt.

Zur rechten Beurtheilung dieses Falles verweisen wir auf eine Regel, die sich in Dr. Walthers „Pastorale“ (S. 339) findet und die mit Gottes Wort erhärtet wird. In den Bann gethan werden können die nicht, die noch nicht nach göttlicher Ordnung von ihrem Irrthum oder von ihrer Sünde fruchtlos überzeugt, ermahnt und gestraft, und noch nicht so als halsstarrige und unverbesserliche Irrgeister oder Sünder offenbar geworden sind. (Matth. 18, 15—17. 2 The. 3, 14. 15. Vgl. Tit. 3, 10. 11.)“ Wir finden, daß die werte Gemeinde in dem Falle A. der Regel des Wortes Gottes: „Höret die Gemeine nicht“ zuwidergehandelt hat. Die Verhandlungen mit Herrn A. zeigen, daß derselbe der brüderlichen Ermahnung, sein Kind in die Gemeindeschule zu schicken, sich nicht verschloß, sondern dem Verlangen, dies sogleich zu thun, mit dem Versprechen entgegentrat, dem nach Verlauf einiger Wochen nachzukommen. Daß die Gemeinde darauf nicht einging, sondern darauf bestand, sein Unrecht sogleich gutzuthun, und Herrn A. für einen Selbstgebannten erklärte, war den Worten Gottes zuwider, da Herr A. sich den Ermahnungen der Gemeinde, das zu thun, was der ganzen Gemeinde und seinem Kinde zu Besten dienen mußte, keineswegs verschloß.

Ferner sagt Dr. Walther, daß man nur dann eine Person für einen Selbstgebannten erklären sollte, „wenn der Grund der über den Betreffenden verhängten Kirchenzucht eine offenbar Todsünde war“. (A. a. O. S. 347.) Wir finden, daß die werte Dreifaltigkeits-Gemeinde nicht unterschieden hat, was nach Gottes Wort eine Sünde ist, in der ein bußfertiger Christ auf keinen Augenblick beharren kann, und dem, was als Durchbrechung einer guten Ordnung betrachtet werden muß. Eine solche gute, löbliche Ordnung ist die Regel der Gemeinde, daß alle Gemeindeglieder gehalten sind, ihre Kinder bis zur Confirmation in die Gemeindeschule zu schicken. Beim Halten auf die Ordnung muß, da es eine gute, aber doch menschliche Ordnung ist, den Verhältnissen, wie auch der Schwachheit Rechnung getragen werden. Die Gemeinde hat aber in dem Fall A., resp. mit dem Ausschluß des Herrn A., mit ihrer Bann-Erklärung, der Regel des Wortes Gottes zuwider gehandelt, da der Grund des über Herrn A. ausgesprochenen Bannurtheils keine offenbar Todsünde, nämlich keine Sünde war, mit der der Glaube nicht bestehen kann.

\subsection{II}

Herr A. hat in seiner Erbitterung darüber, daß er von der Gemeinde als Gebannter erklärt worden war, hinter dem Rücken der Gemeinde und ihr zuwider eine englische Gemeinde ins Leben zu rufen sich entschlossen und demgemäß gehandelt. Das ist unbedingt als sündlich zu bezeichnen, und alle diejenigen, welche in gleicher Gesinnung in diesem Werke seine Mithelfer gewesen sind, haben sich derselben Sünde theilhaftig gemacht. Es war daher auch recht, daß die Gemeinde an diesem Werke keinen Gefallen fand, sondern energisch dagegen Protest erhob. Wir müssen daher auch dem Urtheil der Commission von 1901, dahingehend, daß Herr B. sich durch seine Theilhaberschaft an diesem Werke und da er den Protest der Dreifaltigkeits-Gemeinde nicht gehörig beachtete, dem Vorwurf unbrüderlicher Praxis ausgesetzt habe, voll und ganz beistimmen.

Ferner müssen wir sagen: Wir finden, daß in der neu entstandenen englischen Gemeinde Leute sind, deren bisheriges Verhältniß zur
\quellseite{5}
Dreifaltigkeits-Gemeinde es dieser unmöglich macht, die betreffende englische Gemeinde sofort als Schwester-Gemeinde anzuerkennen. Es sind erstens solche, die zur Zeit der Entstehung der englischen Gemeinde noch Glieder der deutschen Gemeinde waren, wie die Herren S. und D. Die Gemeinde thut recht, von denselben zu verlangen, daß sie nicht nur es als unrecht erkennen, der englischen Gemeinde beigetreten zu sein, ohne ihre bisherige Gemeinde um Entlassung gebeten zu haben, sondern auch bußfertig zu derselben zurückzukehren. Es sind zum andern solche, welche durch Einsendung eines Resignationsschreibens ihren Austritt aus der Gemeinde erklärten, wie z. B. Frau C. und Frau F. Die Gemeinde thut recht, wenn sie von denselben verlangt, daß sie zu ihr zurückkehren und in ordentlicher Weise um Entlassung bitten. Was eine dritte Klasse betrifft, nämlich solche Personen, welche früher mit der deutschen Gemeinde in gliedlicher Verbindung gestanden haben, hernach aber sich von derselben getrennt haben, Wort- und Sacramentsverächter geworden sind und auch nicht mehr von der Gemeinde als zu ihr gehörig behandelt worden sind, wie z. B. die Herren G. und H., so ist unser Urtheil folgendes: Die Gemeinde that recht, von diesen zu fordern, daß sie ihre Gesinnungsänderung und bußfertige Rückkehr zur Kirche ihr zur Kenntniß brächten, ehe sie dieselben wieder als Brüder anerkennen könnte. Doch hat die Gemeinde kein Recht, von denselben zu fordern, daß sie ihre Bußfertigkeit durch Rückkehr in die deutsche Gemeinde beweisen, noch gebietet ihr die Sorge um das Seelenheil derselben, solche Forderung zu stellen; vielmehr sollte gerade die Sorge um deren Seelenheil die Gemeinde bewegen, von dieser Forderung abzusehen, und sollte sie sich freuen, daß die Verirrten sich wieder zu Wort und Sacrament, und zwar in der rechtgläubigen Kirche, halten wollen.

\subsection{III}

Ob das Entlassungsgesuch der Herren J. und R. ein billiges und wohlbegründetes war, entzieht sich unserer Beurtheilung, da die Gemeinde die Gründe dafür nicht angehört und daher auch ihrem Protokoll nicht einverleibt hat. Doch können wir nicht umhin, der lieben Gemeinde zu sagen, daß es nach unserem Dafürhalten ungerecht und unbrüderlich war, die Bittsteller so lange hinzuhalten, ihnen Gehör zu versagen und sie sogar vom heiligen Abendmahl zu suspendiren.

\subsection{IV}

In Bezug auf die ehemalige Christ Church und die jetzige Church of our Savior geben wir folgende Erklärung ab: 1. Wir halten dafür, daß die ehemalige Christ Church und die jetzige Church of our Savior insofern wesentlich ein Ding sind, als die Glieder der ersteren Glieder der letzteren sind. 2. Wir urtheilen, daß die Dreifaltigkeits-Gemeinde die Church of our Savior unter den folgenden Bedingungen anerkennen sollte: a. Die Church of our Savior erklärt ausdrücklich die Art und Weise, wie die Christ Church entstand, nämlich als Oppositionswerk gegen die deutsche Gemeinde, für sündlich; b. die Church of our Savior wird nur solche ehemalige Glieder der Dreifaltigkeits-Gemeinde in ihren Verband aufnehmen, resp. in ihrem Verband behalten, die um eine ordentliche Entlassung nachgesucht und dieselbe erhalten haben.

(gez.) J. Pieper.
„ P. Brand.
„ C. E. Schmidt.

Unsere Meinung war, an diesem Abend Punkt für Punkt mit der Gemeinde zu besprechen, sie von der Richtigkeit unseres Urtheils zu überzeugen, und daß dann der ganze Fall auf Grund und nach Anleitung dieser Sätze zur Erledigung gebracht werden könnte. Wir sind aber den ganzen langen Abend nicht über das erste Stück, den Bannfall, hinausgekommen. Alles Reden und Belehren aus Gottes Wort konnte die Gemeinde nicht bewegen, sich zu unserer Entscheidung zu bekennen und den falschen Bann aufzuheben. Selbstverständlich haben wir der Gemeinde auch mit vielen Worten dies gesagt, daß wir Herrn A. in seinem Verhalten, soweit es sündlich war, durchaus nicht in Schutz nehmen oder auch nur im geringsten entschuldigen wollten, daß sich derselbe nach unserer Ueberzeugung versündigt und die Gemeinde nur ihre Pflicht gethan habe, indem sie sehr ernstlich mit ihm handelte.

Zum zweiten Punkt des Circulars, die Anerkennung der englischen Gemeinde betreffend, sei dies bemerkt: Wir glaubten, als wir nach Cincinnati gingen, es möchte zur Beilegung des Streites dienlich sein, wenn Herr Präses Meyer von der englischen Synode und etwa auch einige Glieder der Church of Our Savior den Verhandlungen beiwohnen könnten. Herr Präses Meyer war zu dem Zweck auch nach Cincinnati gekommen. Aber man weigerte sich ganz entschieden, auch nur Präses Meyer zuzulassen. So konnten wir das gegenseitige Verhältniß der deutschen und englischen Gemeinde nur auf Grund des Materials begutachten, welches uns zu dem Zweck von der deutschen Gemeinde zur Verfügung gestellt wurde. Wir urtheilten demgemäß, wie oben berichtet, daß die englische Gemeinde in sündlicher Weise entstanden sei, sich in verschiedener Beziehung gegen die deutsche versündigt habe und nur unter gewissen Bedingungen anzuerkennen sei.

Unter Punkt 3 des Ultimatums heißt es dann noch: „Schließlich meldete er“ (Herr Präses Riemann ist gemeint) „uns kurz, er habe den Handel dem Allgemeinen Präsidium übergeben, von diesem aber ist bis zu dieser Stunde noch keine officielle Kundgebung eingelaufen bezüglich der von uns erhobenen Klagen und Beschwerden. Offenbar wollte man sich mit einer Untersuchung derselben nicht befassen, nicht die Finger verbrennen.“ Diese Worte haben zu der Auffassung Veranlassung gegeben, als ob der Allgemeine Präses sich überhaupt nicht um die Cincinnati-Angelegenheit gekümmert habe. Die Beschuldigung, der Allgemeine Präses habe in dem Handel seine Pflicht nicht gethan, sich der Gemeinde in Cincinnati in ihrer Noth nicht angenommen, ist gerade auch von der Gemeinde selbst erhoben worden, als wir mit ihr handelten, und Herr Präses Pieper hat nicht unterlassen, die Gemeinde sofort aus ihrem eigenen Protokoll zu überführen, daß ihre Beschuldigung durchaus grundlos sei. Thatsache ist allerdings, daß der Allgemeine Präses absolut keinen Beruf hatte, sich von Amts wegen der Sache anzunehmen, ehe der Präses der englischen Synode bei ihm wider die Gemeinde in Cincinnati klagbar wurde. Jeder, der unsere Synodalordnung kennt, wird sich davon leicht überzeugen können, wenn er hört, daß die ganze Streitsache bis zur Klage des englischen Präses offiziell nur dem Mittleren District und seine Beamten anging. Auch Klagen gegen Districtsbeamte müssen zunächst dem District vorgelegt werden. Doch hat Herr Präses Pieper auf privatem Wege sein Möglichstes gethan zur Heilung des Schadens in Cincinnati, wenn er zu heilen wäre. So schrieb er zunächst an Herrn P. v. Schlichten und an den Vorstand der Gemeinde, als diese ihm ausführliche Mittheilungen über den Streithandel gemacht hatten, er rathe ihnen, die Sache verständigen Männern vorzulegen. P. v. Schlichten theilte dies seiner Gemeinde mit, fügte aber hinzu, wie das Gemeindeprotokoll ausweist,
\quellseite{6}
den Rath des Präses brauche man nicht zu befolgen. Später, als Präses Pieper erfuhr, P. v. Schlichten habe besonderes Vertrauen zu zwei bestimmten Brüdern im Amte, schlug er ihm vor, doch ja diese beiden Brüder zu bitten, nach Cincinnati zu kommen, da er dieselben für überaus geeignete Männer halte. Der Präses erhielt aber den Bescheid, daß Pastor und Gemeinde nicht geneigt seien, auf diesen Vorschlag einzugehen. Endlich bot der Allgemeine Präses es P. v. Schlichten an, daß er selbst nach Cincinnati kommen wolle, mit seinem Rathe zu dienen, daß er hierzu aber von Pastor und Gemeinde aufgefordert werden müsse, da er amtlich keinen Beruf habe, einzugreifen. Auch auf dieses Anerbieten des Präses erfolgte keine zusagende Antwort. In einem Briefe befand sich auch die Bemerkung, daß die Gemeinde alt genug sei und es nicht für ihre Aufgabe halte, um eine Untersuchung zu bitten. —

Nun noch einige Worte über die weitere Entwicklung der Angelegenheit, soweit das Allgemeine Präsidium dabei in Betracht kommt. Bald nach der oben beschriebenen Untersuchung theilten eine Anzahl Glieder der Gemeinde dieser und gleichzeitig dem Allgemeinen Präses mit, daß sie sich jetzt zu dem vorgelegten Gutachten bekennen müßten. Sie wurden darauf von Herrn P. v. Schlichten jun. vom Abendmahl suspendirt, und man drohte ihnen mit einem Zuchtverfahren. Präses Pieper bat hierauf die Gemeinde, ihm Gelegenheit zu einer nochmaligen Verhandlung zu geben. Die Gemeinde antwortete, daß sie vorläufig nicht weiter handeln möchte, da sie andere Sachen zu thun habe. Präses Pieper richtete dann am 29. December folgendes Schreiben an die Gemeinde:

St. Louis, Mo., den 29. December 1902.

\section[An die Dreifaltigkeits-Gemeinde]{An die ev.-luth. Dreifaltigkeits-Gemeinde zu Cincinnati, O.}

In Christo herzlich geliebte Brüder!

Ihr habt mich in einem Schreiben, das am 6. December in meine Hände gelangt ist, benachrichtigt, daß Ihr vorläufig nicht weiter in Sachen des von Euch ausgesandten Circulars mit den Beamten der Synode zu verhandeln wünscht. So gerne nun die Beamten der Synode sich in die Gemeinden schicken, so muß ich Euch doch daran erinnern, daß die Sache, um die es sich handelt, keinen langen Aufschub duldet. Wir, die Beamten der Synode, müssen der Kirche, an die Ihr Euer Circular gesendet habt, möglichst bald darüber Bericht erstatten, was wir bei Euch gefunden haben. Sodann ist eine weitere Behandlung der Angelegenheit auch um Eurer selbst willen dringend nöthig, weil Ihr — das geht aus Eurem letzten Schreiben deutlich hervor — noch nicht zur rechten Klarheit aus und nach Gottes Wort in der streitigen Sache gekommen seid. So bitte ich Euch denn um Christi willen, daß Ihr Euch nicht länger weigert, in der Sache alsbald mit den Synodalbeamten weiter zu handeln. Ich stelle Euch meine Antwort gedruckt zu, weil Ihr mir mittheilt, Ihr hättet keine Zeit, die Angelegenheit in der Gemeindeversammlung aufzunehmen. So hat jedes Gemeindeglied Gelegenheit, meine freundliche Erinnerung zu lesen.

Ich gebe noch einmal die Punkte an, in denen es die Gemeinde versehen hat.

1. Die Gemeinde hat geirrt in dem A.’schen Bannfall. Wollt Folgendes sorgfältig und vor Gott bei Euch erwägen: Wenn eine christliche Gemeinde in Kirchenzuchtsfällen auf dem rechten Wege bleiben will, so muß sie zwischen Sünden und Sünden wohl unterscheiden. Es gibt Sünden, bei denen der Glaube an Christum keinen Augenblick bestehen kann, z. B. ehebrechen, stehlen. Bei solchen Sünden muß die Gemeinde darauf bestehen, daß sie sofort erkannt und abgethan werden. Einem Dieb sagt die Gemeinde mit Recht: du mußt das Stehlen als Sünde erkennen und sofort lassen, oder du kannst nicht in der Gemeinschaft der Christen bleiben. Hier ist kein Frist bis Neujahr oder Ostern zu geben. Daneben gibt es Sünden bei Christen, die den Glauben nicht sofort austreiben, Sünden, deren Ursache Schwäche in der Erkenntniß, oder, bei guter Erkenntniß, zeitweiliges Schwachwerden im Glauben zu Grunde liegen kann, z. B. Zugehörigkeit zu sündlichen Arbeiter-Unions, Trusts etc. In diesen Fällen darf die Gemeinde nicht sagen: Wenn du nicht sofort austrittst, betrachten wir dich als einen Gebannten oder Selbstgebannten. Kurz, jede Gemeinde, die bei der Kirchenzucht auf rechtem Wege bleiben will, muß sorgfältig unterscheiden zwischen Sünden, bei denen der Glaube nicht bestehen kann (Todssünden), und zwischen Sünden, bei denen das Bestehen des Glaubens an Christum noch möglich ist (Schwachheitssünden). Bei Sünden, bei denen der Glaube nicht bestehen kann, verhängt die Gemeinde einen rechten Bann, wenn die Sünde auf ihre Mahnung hin nicht erkannt und sofortige Abthuung versprochen wird. Bei Sünden hingegen, bei denen der Glaube noch bestehen kann, verhängt die Gemeinde einen falschen Bann, wenn sie jemand für einen Gebannten oder Selbstgebannten erklärt, der nicht sofort die Sünde erkennt und abthut. Wenden wir dies nun an auf den A.’schen Fall. Herr A. wollte sein Kind ehe aus der Gemeindeschule nehmen und confirmirt haben, als es die Gemeinde für recht hielt. Die Gemeinde ging nicht darauf ein. Herr A. nahm darauf sein Kind aus der Gemeindeschule und sandte es in die public school. Die Gemeinde verhandelte mit ihm. A. versprach, sein Kind nächsten September und, nach weiteren Verhandlungen nach Neujahr wieder in die Gemeindeschule senden zu wollen. Diese Verhandlungen fanden etwa im October statt. A. bat um diese Frist seiner schwierigen häuslichen Verhältnisse wegen. Die Gemeinde verweigerte jegliche Frist und forderte als Zeichen der Bußfertigkeit sofortige Sendung des Kindes in die Gemeindeschule. Als hierauf A. langsam die Gemeindeversammlung verließ, indem er Worte sagte wie diese: „Er sehe, die Gemeinde wolle ihn los werden“ etc., erklärte ihn die Gemeinde für einen Selbstgebannten und forderte in der öffentlichen Bannserklärung auf, A. als einen Heiden und Zöllner zu behandeln. Das ist in Kürze der Handel mit Herrn A. Wie ist er zu beurtheilen? Die Gemeinde hat entschieden verkehrt gehandelt. Die Gemeinde konnte wohl überzeugt sein, daß Herr A. sich versündigte, indem er sein Kind zu früh aus der Gemeindeschule nehmen wollte und dadurch für das geistliche Wohl seines Kindes nicht so sorgte, wie er es hätte thun sollen etc. Aber die Gemeinde konnte nicht überzeugt sein, daß in A.’s Herzen jedes Fünklein Glaube erloschen war. Darum hat die Gemeinde einen falschen Bann verhängt, weil eine christliche Gemeinde nach Gottes Wort nur dann den Bann verhängen darf, wenn die Sünde der Art ist, daß sie unbedingt den Glauben ausschließt und ohne Frucht aus Gottes Wort gestraft wurde. So schreibt der Apostel Paulus 1 Cor. 5, 13.: „Thut von euch selbst hinaus, wer da böse ist.“ Dr. Walther sagt daher in seinem „Pastorale“ (S. 338), daß nur ein solcher in den Bann gethan werden könne, der „als ein unverbesserlicher Unchrist offenbar geworden ist“. Eine Gemeinde darf daher nie den Bann verhängen, wenn noch die Annahme möglich ist, daß in dem Sündigenden noch ein Fünklein Glaube vorhanden sei. Wird der Bann verhängt, trotzdem noch ein Fünklein Glaube in dem Sündigenden vorhanden ist, so trifft der Bann den Gebannten nicht, sondern fällt als eine schwere Sünde auf die Gemeinde. Dieser schweren Sünde hat sich die Gemeinde in Cincinnati, wenn auch in guter Meinung, schuldig gemacht. Diese Sünde muß sie abthun.

\quellseite{7}

2. Die Gemeinde in Cincinnati irrt ferner darin, daß sie es als eine unberufene Einmischung in ihre Gemeindeangelegenheiten ansieht, wenn Synodalbeamte die Rechtmäßigkeit des von ihr verhängten Bannes untersuchen. Die Sache steht so: jede christliche Gemeinde, auch wenn sie lediglich zu keiner Synode gehört, ist verpflichtet, jeder andern Gemeinde Einsicht in ihr Kirchenzuchtverfahren zu gewähren, wenn die Rechtmäßigkeit desselben angefochten wird. Dies fordert die Liebe und die Rücksicht auf die Gewissen, welche in Bezug auf das kirchenzuchtliche Verfahren beunruhigt sind. Will eine Gemeinde ihr Kirchenzuchtverfahren nicht prüfen lassen, so setzt sie sich dem Verdacht aus, daß sie keine gute Sache habe, und die Gemeinde kann nicht verlangen, daß der von ihr verhängte Bann von andern respektirt werde. Deshalb ermahnt Dr. Walther in seinem „Pastorale“ (S. 351) die Gemeinden, über alle Kirchenzuchtsfälle genau Protokoll zu führen, und setzt hinzu: „Die Gemeinde sollte allezeit aus ihrem Protokoll die Richtigkeit ihres Verfahrens in jedem vorgekommenen Bannfall nachweisen können, da ohne diesen Nachweis andere Gemeinden nicht in der Lage sind, den Bann auf alle Fälle respektiren zu können oder zu müssen.“ Vollends sollte eine Gemeinde, die Glied der Synode ist, sich nicht weigern, in einem Klagefalle einen verhängten Bann durch die Synodalbeamten untersuchen zu lassen. Die gegenseitige Ueberwachung bei Bannfällen ist gerade auch ein Zweck der Synodalverbindung. Es heißt in unserem „Synodalhandbuch“ (S. 148): „Die Synode bittet und ermahnt alle Gemeinden und Prediger recht herzlich, bei einem vorkommenden Bannfalle nicht allein handeln zu wollen, sondern die Hülfe und den Rath entweder des betreffenden Districtspräses oder anderer älterer und erfahrener Amtsbrüder in Anspruch zu nehmen.“ Ferner wird im „Synodalhandbuch“ (S. 145) den Visitatoren zur Pflicht gemacht: „Die Visitatoren sollen ebenfalls bei ihren Visitationen regelmäßig Nachfrage halten, ob Bannfälle vorgekommen seien, und Einsicht in die betreffenden Protokolle nehmen.“ Diese Synodalordnung hat auch die Gemeinde in Cincinnati zu der ihrigen gemacht. Es ist daher sehr verkehrt, wenn die Gemeinde die Untersuchung des A.’schen Bannfalles durch die Synodalbeamten als eine unberechtigte Einmischung in Gemeindeangelegenheiten ansehen will. — Auf die Frage, welche nebenbei aufgeworfen ist, ob eine christliche Ortsgemeinde in Bann sache irren könne, da ihr doch der Bann übergeben sei, ist zu antworten: eine christliche Gemeinde, wie auch der einzelne Christ, kann nicht irren, wenn sie nach Gottes Wort handelt. Sobald sie aber nicht nach Gottes Wort, sondern nach dem eigenen Kopfe handeln, sind sie dem Irrthum verfallen. Daß auch Gemeinden in Irrthum geraten können, sehen wir aus Gal. 3, 1., wo der Apostel Paulus den galatischen Gemeinden zuruft: „O ihr unverständigen Galater, wer hat euch bezaubert, daß ihr der Wahrheit nicht gehorchet?“

3. In der Gemeinde zu Cincinnati ist ein großer Eifer für die Gemeindeschule. In Bezug auf den Eifer für die Gemeindeschule ist die Gemeinde mit der ganzen Synode einig. Die Synode wird auch fernerhin mit ganzem Ernst auf das kirchliche Institut der Gemeindeschule dringen. Aber die Gemeinde ist entschieden auf den Irrweg gekommen in Bezug auf die Art und Weise, wie man die Gemeindeglieder anhält, ihre Kinder in die Gemeindeschule zu senden. Der selige Dr. Walther wäre, wenn er in der Gemeinde in Cincinnati gelebt hätte, in Gefahr gekommen, von der Gemeinde gebannt zu werden. Ich will ein Vorkommniß aus unserer Gemeinde berichten. Einige Jahre vor dem Tode Dr. Walthers wurde in unserer Gemeindeversammlung ein Zusatz zu der Gemeindeordnung beantragt, des Inhalts, daß jedes Gemeindeglied verpflichtet sein solle, seine Kinder in die Gemeindeschule zu senden. Sofort erhob sich Dr. Walther und legte Protest gegen den Vorschlag ein, weil die Gemeinde kein Recht habe, einen solchen Beschluß zu fassen. Er führte dann weiter aus: eine Gemeinde kann ihren Gliedern nur das gebieten, was Gottes Wort gebietet. Gottes Wort aber gebietet den Eltern hinsichtlich ihrer Kinder, daß sie dieselben auferziehen in der Zucht und Vermahnung zum Herrn. Die christliche Erziehung der Kinder muß daher die Gemeinde von jedem Gemeindeglied fordern und nötigenfalls öffentliche Zucht anwenden. Sie kann es nicht in das Belieben der Gemeindeglieder stellen, ob sie ihre Kinder christlich erziehen wollen oder nicht. Was aber das kirchliche Institut der Gemeindeschule betrifft, so ist diese nur ein Mittel, wenn auch ein ganz vorzügliches, das wir in christlicher Freiheit und Weisheit anwenden, um jenes göttliche Gebot der christlichen Kindererziehung in Ausführung bringen zu helfen. Wir können es den Eltern nicht verbieten, daß sie ihre Kinder selbst unterrichten, oder sich einen Hauslehrer halten etc. In dieser Beziehung haben wir es allerdings in das Belieben der Gemeindeglieder zu stellen, ob sie die Gemeindeschule benutzen wollen oder nicht. Deshalb enthalten die echt-synodischen Gemeindeordnungen, die einen Paragraphen über Gemeindeschulen haben, etwa diese Bestimmung: „Alle Gemeindeglieder sind gehalten, ihre Kinder in die Gemeindeschule zu senden, oder sonst für den christlichen Unterricht ihrer Kinder Sorge zu tragen.“ Ob letzteres geschieht, hat die christliche Gemeinde allerdings zu überwachen. Aber sie kann nicht schlechthin festsetzen: „Alle Kinder müssen in die Gemeindeschule gehen.“ Sie kann daher auch nicht im einzelnen Falle schlechthin beschließen: „Du sendest dein Kind sofort in die Gemeindeschule, oder du wirst in Kirchenzucht genommen und von der Gemeinde ausgeschlossen.“ Die Gemeinde muß sich mit ihrer Forderung auf die Forderung der christlichen Kindererziehung beschränken. Das ist göttliches Gebot. Aber gehen bei dieser Weise nicht die Gemeindeschulen ein? Durchaus nicht, wie die Geschichte unserer Synode beweist. Gerade bei dieser Weise werden die Gemeindeschulen gefüllt. Wie fangen wir das an? So, daß wir immerfort die Gewissen in Bezug auf das göttliche Gebot: „Ziehet eure Kinder auf in der Zucht und Vermahnung zum Herrn“ schärfen. Dabei weisen wir hin auf die Herrlichkeit der Gemeindeschule, in welcher die Kinder viel mehr lernen als in allen public schools, high schools und Universitäten, nämlich Gottes Wort, wie man selig wird, die rechte christliche Weltanschauung etc. Wir weisen auch hin auf die Gefahren, welche für den Glauben der ChristenKinder im Besuch der öffentlichen Schulen liegen. Wenn wir so die Herrlichkeit des kirchlichen Instituts der Gemeindeschule hervorheben und dabei, wie gesagt, die Eltern immer an das göttliche Gebot, ihre Kinder in der Zucht und Vermahnung zum Herrn zu erziehen, erinnern, so bewirken wir, daß christliche Eltern mit Freuden nach der Gemeindeschule greifen und Gott von Herzen dafür danken. Wo es einmal fehlen will, da hilft immer wieder die freundliche Erinnerung und Ermahnung. Ich war von Herzen erschrocken, als ich in Eurer Gemeindeversammlung etwa die Worte hörte: „Wir könnten unsere Lehrer entlassen und die Schule schließen, wenn wir niemand mehr wegen Nichtsendung der Kinder in die Gemeindeschule in den Bann thun dürften.“ Darnach würde die Furcht vor dem Bann die Gemeindeschule zusammenhalten. Gott sei Dank, daß es euch in Cincinnati nicht so steht. Aber jene Rede zeigt eine verkehrte, gesetzliche Stellung in Sachen der Gemeindeschule. Eure Gemeindeschule wird gedeihen, wie die Gemeindeschulen in der ganzen Synode, wenn Ihr die Schul sache behandelt, wie Dr. Walther und die ganze Synode.

Ich habe Euch mit vielen Worten geschrieben, weil ich sehe, daß die Gemeinde in der A.’schen Bannsache einen falschen Weg betreten hat und sich weigert, weitere Belehrung anzunehmen. Dies müßte zum
\quellseite{8}
Verderben der Gemeinde ausschlagen. Ich bitte Euch daher herzlich und um Christi willen: zieht Euren gefaßten Beschluß in Wiedererwägung und erklärt Euch bereit, die Sache weiter mit den Synodalbeamten nach Gottes Wort zu verhandeln und zu einem gottgefälligen Abschluß zu bringen.

\begin{unterschrift}
Mit brüderlichem Gruß\\
F. Pieper.
\end{unterschrift}

Aus Veranlassung dieses Schreibens kam es noch einmal zu einer mündlichen Verhandlung mit der Gemeinde, bei welcher Gelegenheit Herr Dr. Gräbner der Begleiter des Präses war. Als jedoch auch diese Bemühung fehlschlug und die Gemeinde ausdrücklich brieflich erklärt hatte, sie müsse bei ihrer Stellung beharren, glaubte der Präses, an seinem Theil alles gethan zu haben, und es blieb ihm jetzt nur noch übrig, ordnungsmäßig dem Mittleren District Anzeige von der Sachlage in Cincinnati zu machen, damit derselbe auf disciplinarischem Wege mit der dortigen Gemeinde weiter handle.

\begin{unterschrift}
C. E. Schmidt.
\end{unterschrift}
